% Options for packages loaded elsewhere
\PassOptionsToPackage{unicode}{hyperref}
\PassOptionsToPackage{hyphens}{url}
\PassOptionsToPackage{dvipsnames,svgnames,x11names}{xcolor}
%
\documentclass[
  12pt]{article}

\usepackage{amsmath,amssymb}
\usepackage{iftex}
\ifPDFTeX
  \usepackage[T1]{fontenc}
  \usepackage[utf8]{inputenc}
  \usepackage{textcomp} % provide euro and other symbols
\else % if luatex or xetex
  \usepackage{unicode-math}
  \defaultfontfeatures{Scale=MatchLowercase}
  \defaultfontfeatures[\rmfamily]{Ligatures=TeX,Scale=1}
\fi
\usepackage{lmodern}
\ifPDFTeX\else  
    % xetex/luatex font selection
\fi
% Use upquote if available, for straight quotes in verbatim environments
\IfFileExists{upquote.sty}{\usepackage{upquote}}{}
\IfFileExists{microtype.sty}{% use microtype if available
  \usepackage[]{microtype}
  \UseMicrotypeSet[protrusion]{basicmath} % disable protrusion for tt fonts
}{}
\makeatletter
\@ifundefined{KOMAClassName}{% if non-KOMA class
  \IfFileExists{parskip.sty}{%
    \usepackage{parskip}
  }{% else
    \setlength{\parindent}{0pt}
    \setlength{\parskip}{6pt plus 2pt minus 1pt}}
}{% if KOMA class
  \KOMAoptions{parskip=half}}
\makeatother
\usepackage{xcolor}
\setlength{\emergencystretch}{3em} % prevent overfull lines
\setcounter{secnumdepth}{5}
% Make \paragraph and \subparagraph free-standing
\makeatletter
\ifx\paragraph\undefined\else
  \let\oldparagraph\paragraph
  \renewcommand{\paragraph}{
    \@ifstar
      \xxxParagraphStar
      \xxxParagraphNoStar
  }
  \newcommand{\xxxParagraphStar}[1]{\oldparagraph*{#1}\mbox{}}
  \newcommand{\xxxParagraphNoStar}[1]{\oldparagraph{#1}\mbox{}}
\fi
\ifx\subparagraph\undefined\else
  \let\oldsubparagraph\subparagraph
  \renewcommand{\subparagraph}{
    \@ifstar
      \xxxSubParagraphStar
      \xxxSubParagraphNoStar
  }
  \newcommand{\xxxSubParagraphStar}[1]{\oldsubparagraph*{#1}\mbox{}}
  \newcommand{\xxxSubParagraphNoStar}[1]{\oldsubparagraph{#1}\mbox{}}
\fi
\makeatother

\usepackage{color}
\usepackage{fancyvrb}
\newcommand{\VerbBar}{|}
\newcommand{\VERB}{\Verb[commandchars=\\\{\}]}
\DefineVerbatimEnvironment{Highlighting}{Verbatim}{commandchars=\\\{\}}
% Add ',fontsize=\small' for more characters per line
\usepackage{framed}
\definecolor{shadecolor}{RGB}{241,243,245}
\newenvironment{Shaded}{\begin{snugshade}}{\end{snugshade}}
\newcommand{\AlertTok}[1]{\textcolor[rgb]{0.68,0.00,0.00}{#1}}
\newcommand{\AnnotationTok}[1]{\textcolor[rgb]{0.37,0.37,0.37}{#1}}
\newcommand{\AttributeTok}[1]{\textcolor[rgb]{0.40,0.45,0.13}{#1}}
\newcommand{\BaseNTok}[1]{\textcolor[rgb]{0.68,0.00,0.00}{#1}}
\newcommand{\BuiltInTok}[1]{\textcolor[rgb]{0.00,0.23,0.31}{#1}}
\newcommand{\CharTok}[1]{\textcolor[rgb]{0.13,0.47,0.30}{#1}}
\newcommand{\CommentTok}[1]{\textcolor[rgb]{0.37,0.37,0.37}{#1}}
\newcommand{\CommentVarTok}[1]{\textcolor[rgb]{0.37,0.37,0.37}{\textit{#1}}}
\newcommand{\ConstantTok}[1]{\textcolor[rgb]{0.56,0.35,0.01}{#1}}
\newcommand{\ControlFlowTok}[1]{\textcolor[rgb]{0.00,0.23,0.31}{\textbf{#1}}}
\newcommand{\DataTypeTok}[1]{\textcolor[rgb]{0.68,0.00,0.00}{#1}}
\newcommand{\DecValTok}[1]{\textcolor[rgb]{0.68,0.00,0.00}{#1}}
\newcommand{\DocumentationTok}[1]{\textcolor[rgb]{0.37,0.37,0.37}{\textit{#1}}}
\newcommand{\ErrorTok}[1]{\textcolor[rgb]{0.68,0.00,0.00}{#1}}
\newcommand{\ExtensionTok}[1]{\textcolor[rgb]{0.00,0.23,0.31}{#1}}
\newcommand{\FloatTok}[1]{\textcolor[rgb]{0.68,0.00,0.00}{#1}}
\newcommand{\FunctionTok}[1]{\textcolor[rgb]{0.28,0.35,0.67}{#1}}
\newcommand{\ImportTok}[1]{\textcolor[rgb]{0.00,0.46,0.62}{#1}}
\newcommand{\InformationTok}[1]{\textcolor[rgb]{0.37,0.37,0.37}{#1}}
\newcommand{\KeywordTok}[1]{\textcolor[rgb]{0.00,0.23,0.31}{\textbf{#1}}}
\newcommand{\NormalTok}[1]{\textcolor[rgb]{0.00,0.23,0.31}{#1}}
\newcommand{\OperatorTok}[1]{\textcolor[rgb]{0.37,0.37,0.37}{#1}}
\newcommand{\OtherTok}[1]{\textcolor[rgb]{0.00,0.23,0.31}{#1}}
\newcommand{\PreprocessorTok}[1]{\textcolor[rgb]{0.68,0.00,0.00}{#1}}
\newcommand{\RegionMarkerTok}[1]{\textcolor[rgb]{0.00,0.23,0.31}{#1}}
\newcommand{\SpecialCharTok}[1]{\textcolor[rgb]{0.37,0.37,0.37}{#1}}
\newcommand{\SpecialStringTok}[1]{\textcolor[rgb]{0.13,0.47,0.30}{#1}}
\newcommand{\StringTok}[1]{\textcolor[rgb]{0.13,0.47,0.30}{#1}}
\newcommand{\VariableTok}[1]{\textcolor[rgb]{0.07,0.07,0.07}{#1}}
\newcommand{\VerbatimStringTok}[1]{\textcolor[rgb]{0.13,0.47,0.30}{#1}}
\newcommand{\WarningTok}[1]{\textcolor[rgb]{0.37,0.37,0.37}{\textit{#1}}}

\providecommand{\tightlist}{%
  \setlength{\itemsep}{0pt}\setlength{\parskip}{0pt}}\usepackage{longtable,booktabs,array}
\usepackage{calc} % for calculating minipage widths
% Correct order of tables after \paragraph or \subparagraph
\usepackage{etoolbox}
\makeatletter
\patchcmd\longtable{\par}{\if@noskipsec\mbox{}\fi\par}{}{}
\makeatother
% Allow footnotes in longtable head/foot
\IfFileExists{footnotehyper.sty}{\usepackage{footnotehyper}}{\usepackage{footnote}}
\makesavenoteenv{longtable}
\usepackage{graphicx}
\makeatletter
\def\maxwidth{\ifdim\Gin@nat@width>\linewidth\linewidth\else\Gin@nat@width\fi}
\def\maxheight{\ifdim\Gin@nat@height>\textheight\textheight\else\Gin@nat@height\fi}
\makeatother
% Scale images if necessary, so that they will not overflow the page
% margins by default, and it is still possible to overwrite the defaults
% using explicit options in \includegraphics[width, height, ...]{}
\setkeys{Gin}{width=\maxwidth,height=\maxheight,keepaspectratio}
% Set default figure placement to htbp
\makeatletter
\def\fps@figure{htbp}
\makeatother

\addtolength{\oddsidemargin}{-.5in}%
\addtolength{\evensidemargin}{-1in}%
\addtolength{\textwidth}{1in}%
\addtolength{\textheight}{1.7in}%
\addtolength{\topmargin}{-1in}%
\usepackage{booktabs}
\usepackage{longtable}
\usepackage{array}
\usepackage{multirow}
\usepackage{wrapfig}
\usepackage{float}
\usepackage{colortbl}
\usepackage{pdflscape}
\usepackage{tabu}
\usepackage{threeparttable}
\usepackage{threeparttablex}
\usepackage[normalem]{ulem}
\usepackage{makecell}
\usepackage{xcolor}
\makeatletter
\@ifpackageloaded{caption}{}{\usepackage{caption}}
\AtBeginDocument{%
\ifdefined\contentsname
  \renewcommand*\contentsname{Table of contents}
\else
  \newcommand\contentsname{Table of contents}
\fi
\ifdefined\listfigurename
  \renewcommand*\listfigurename{List of Figures}
\else
  \newcommand\listfigurename{List of Figures}
\fi
\ifdefined\listtablename
  \renewcommand*\listtablename{List of Tables}
\else
  \newcommand\listtablename{List of Tables}
\fi
\ifdefined\figurename
  \renewcommand*\figurename{Figure}
\else
  \newcommand\figurename{Figure}
\fi
\ifdefined\tablename
  \renewcommand*\tablename{Table}
\else
  \newcommand\tablename{Table}
\fi
}
\@ifpackageloaded{float}{}{\usepackage{float}}
\floatstyle{ruled}
\@ifundefined{c@chapter}{\newfloat{codelisting}{h}{lop}}{\newfloat{codelisting}{h}{lop}[chapter]}
\floatname{codelisting}{Listing}
\newcommand*\listoflistings{\listof{codelisting}{List of Listings}}
\makeatother
\makeatletter
\makeatother
\makeatletter
\@ifpackageloaded{caption}{}{\usepackage{caption}}
\@ifpackageloaded{subcaption}{}{\usepackage{subcaption}}
\makeatother

\ifLuaTeX
  \usepackage{selnolig}  % disable illegal ligatures
\fi
\usepackage[]{natbib}
\bibliographystyle{agsm}
\usepackage{bookmark}

\IfFileExists{xurl.sty}{\usepackage{xurl}}{} % add URL line breaks if available
\urlstyle{same} % disable monospaced font for URLs
\hypersetup{
  pdftitle={Detroit Transportation Survey: Multi-Stage Area Probability Sample Design},
  pdfauthor={Group Name},
  colorlinks=true,
  linkcolor={blue},
  filecolor={Maroon},
  citecolor={Blue},
  urlcolor={Blue},
  pdfcreator={LaTeX via pandoc}}



\begin{document}


\def\spacingset#1{\renewcommand{\baselinestretch}%
{#1}\small\normalsize} \spacingset{1}


%%%%%%%%%%%%%%%%%%%%%%%%%%%%%%%%%%%%%%%%%%%%%%%%%%%%%%%%%%%%%%%%%%%%%%%%%%%%%%

\date{March 6, 2026}
\title{\bf Detroit Transportation Survey: Multi-Stage Area Probability
Sample Design}
\author{
Group Name\\
University of Michigan\\
}
\maketitle

\bigskip
\bigskip
\begin{abstract}

\end{abstract}


\newpage
\spacingset{1.9} % DON'T change the spacing!

\renewcommand*\contentsname{Table of contents}
{
\hypersetup{linkcolor=}
\setcounter{tocdepth}{3}
\tableofcontents
}

\subsection{\texorpdfstring{These are the packages we will be using. The
\href{https://walker-data.com/tidycensus/}{tidycensus} is the only
package that requires an api
key}{These are the packages we will be using. The tidycensus is the only package that requires an api key}}\label{these-are-the-packages-we-will-be-using.-the-tidycensus-is-the-only-package-that-requires-an-api-key}

\subsubsection{Once you request your api key, you will need to save it
locally and not share it in any of the scripts that will go on github.
We may also use the stadia package for
plots}\label{once-you-request-your-api-key-you-will-need-to-save-it-locally-and-not-share-it-in-any-of-the-scripts-that-will-go-on-github.-we-may-also-use-the-stadia-package-for-plots}

\begin{Shaded}
\begin{Highlighting}[]
\CommentTok{\# paths may vary}
\NormalTok{api\_keys }\OtherTok{\textless{}{-}} \FunctionTok{read\_csv}\NormalTok{(}\StringTok{"D:/repos/api{-}keys.csv"}\NormalTok{, }\AttributeTok{show\_col\_types =} \ConstantTok{FALSE}\NormalTok{)}
\end{Highlighting}
\end{Shaded}

\subsection{The data for this project can be read in through through
github}\label{the-data-for-this-project-can-be-read-in-through-through-github}

\begin{Shaded}
\begin{Highlighting}[]
\CommentTok{\# change as needed path to the github repository}
\NormalTok{kevin\_input\_path }\OtherTok{\textless{}{-}} \StringTok{"D:/repos/eight{-}mile{-}sample/inputs/"}
\CommentTok{\# tavril\_input\_path \textless{}{-} }

\NormalTok{census\_tract }\OtherTok{\textless{}{-}} \StringTok{"Census\_Tract\_FIPS\_Code\_Detroit\_MI.csv"}
\NormalTok{detroit\_mi }\OtherTok{\textless{}{-}} \StringTok{"Detroit\_MI.xlsx"}


\CommentTok{\# read in census tracks}
\NormalTok{tract\_fips }\OtherTok{\textless{}{-}} \FunctionTok{read\_csv}\NormalTok{(}
  \FunctionTok{str\_c}\NormalTok{(kevin\_input\_path, }\CommentTok{\# change this as needed}
\NormalTok{        census\_tract), }\AttributeTok{col\_names =} \StringTok{"GEOID"}\NormalTok{,}
  \AttributeTok{col\_types =} \FunctionTok{cols}\NormalTok{(}\AttributeTok{GEOID =} \FunctionTok{col\_character}\NormalTok{()),}
  \AttributeTok{skip =} \DecValTok{1}
\NormalTok{)}

\CommentTok{\# read in detroit tracks}
\NormalTok{tract }\OtherTok{\textless{}{-}} \FunctionTok{read\_excel}\NormalTok{(}
  \FunctionTok{str\_c}\NormalTok{(kevin\_input\_path,}
\NormalTok{        detroit\_mi), }\AttributeTok{sheet =} \StringTok{"Tract"}\NormalTok{) }\SpecialCharTok{|\textgreater{}}
  \FunctionTok{mutate}\NormalTok{(}\AttributeTok{Tract =} \FunctionTok{as.character}\NormalTok{(Tract)) }\SpecialCharTok{|\textgreater{}}
  \FunctionTok{semi\_join}\NormalTok{(tract\_fips, }\AttributeTok{by =} \FunctionTok{c}\NormalTok{(}\StringTok{"Tract"} \OtherTok{=} \StringTok{"GEOID"}\NormalTok{))}

\NormalTok{bg }\OtherTok{\textless{}{-}} \FunctionTok{read\_excel}\NormalTok{(}
   \FunctionTok{str\_c}\NormalTok{(kevin\_input\_path,}
\NormalTok{        detroit\_mi), }\AttributeTok{sheet =} \StringTok{"BlockGroup"}\NormalTok{) }\SpecialCharTok{|\textgreater{}}
  \FunctionTok{mutate}\NormalTok{(}
    \AttributeTok{BlockGroup =} \FunctionTok{as.character}\NormalTok{(BlockGroup),}
    \AttributeTok{Tract =} \FunctionTok{substr}\NormalTok{(BlockGroup, }\DecValTok{1}\NormalTok{, }\DecValTok{11}\NormalTok{)}
\NormalTok{  ) }\SpecialCharTok{|\textgreater{}}
  \FunctionTok{semi\_join}\NormalTok{(tract\_fips, }\AttributeTok{by =} \FunctionTok{c}\NormalTok{(}\StringTok{"Tract"} \OtherTok{=} \StringTok{"GEOID"}\NormalTok{))}

\FunctionTok{cat}\NormalTok{(}\StringTok{"Detroit tracts:"}\NormalTok{, }\FunctionTok{nrow}\NormalTok{(tract), }\StringTok{"}\SpecialCharTok{\textbackslash{}n}\StringTok{"}\NormalTok{)}
\end{Highlighting}
\end{Shaded}

\begin{verbatim}
Detroit tracts: 300 
\end{verbatim}

\begin{Shaded}
\begin{Highlighting}[]
\FunctionTok{cat}\NormalTok{(}\StringTok{"Detroit block groups:"}\NormalTok{, }\FunctionTok{nrow}\NormalTok{(bg), }\StringTok{"}\SpecialCharTok{\textbackslash{}n}\StringTok{"}\NormalTok{)}
\end{Highlighting}
\end{Shaded}

\begin{verbatim}
Detroit block groups: 651 
\end{verbatim}

\newpage

Introduction

This report will capture the the design of a multi-stage area
probability sample for a face-to-face survey of Detroit, Michigan
residents. The survey examines transportation security and
transportation use, with key outcome variables including

\begin{itemize}
\item
  the proportion of residents in households without a vehicle
\item
  the proportion whose main transportation is public transit, the
  proportion who skip destinations due to transportation problems
\item
  the proportion who sought transportation help from friends, family, or
  neighbors.
\end{itemize}

The sample uses a three-stage design:

\begin{enumerate}
\def\labelenumi{\arabic{enumi}.}
\tightlist
\item
  Stage 1: Census tracts as primary sampling units (PSUs),
\item
  Stage 2: Block groups (BGs) as secondary sampling units (SSUs), and
\item
  Stage 2: Persons as elements.
\end{enumerate}

\subsection{Design Goals}\label{design-goals}

The design goals are to yield 1,000 total respondents (see table below)
while achieving a self-weighting sample within each age group and equal
interviewer workload in each sample PSU. These dual objectives are
accomplished through the composite measure of size (MOS) technique
described in VDK.

\subsection{Deliverable Requiremetns}\label{deliverable-requiremetns}

The report will be structured for a dual audience: managers seeking a
high-level understanding of the sample design and its rationale, and
survey statisticians requiring the mathematical details, selection
probabilities, and variance estimation guidance.

Sample Design

\section{Design Goals}\label{design-goals-1}

The sample must satisfy four requirements simultaneously:

\begin{enumerate}
\def\labelenumi{\arabic{enumi}.}
\tightlist
\item
  \textbf{Total sample size}: 1,000 respondents split equally (250 per
  group) across four age domains.
\item
  \textbf{Self-weighting within domains}: Every person in a given age
  group has the same overall selection probability, i.e.,
  \(\pi_k(d) = f_d\) for domain \(d\).
\item
  \textbf{Equal workload per PSU}: The total expected number of selected
  persons is the same across all 40 sample tracts.
\item
  \textbf{Multi-stage selection}: 40 tracts (PSUs), 1 BG per tract
  (SSU), and persons selected within BGs at domain-specific rates.
\end{enumerate}

\section{Target Population}\label{target-population}

The target population consists of all non-institutionalized persons aged
18 and older residing in the City of Detroit, Michigan. According to the
2020 Decennial Census Demographic and Housing Characteristics (DHC)
Summary File, Detroit contains approximately 480,131 adults (18+) across
270 eligible census tracts.

\section{Area Sampling Frame}\label{sec-frame}

\subsection{Data Sources}\label{data-sources}

The sampling frame is constructed from two files:

\begin{itemize}
\item
  \textbf{Census Tract FIPS Code Detroit MI.csv}: A list of 280 FIPS
  codes identifying Census tracts within the City of Detroit, obtained
  from the City of Detroit Open Data Portal.
\item
  \textbf{Detroit MI.xlsx}: Census tract- and block group-level
  aggregate data from the 2020 Decennial Census DHC Summary File,
  downloaded via the \texttt{tidycensus} R package for Wayne County, MI
  and merged with the Detroit tract FIPS codes. Includes housing unit
  counts (\texttt{H1\_001N}), total person counts, and detailed person
  counts by Sex \(\times\) Age (\texttt{P12} table).
\end{itemize}

\subsection{Data cleaning (still in
development)}\label{data-cleaning-still-in-development}

\begin{Shaded}
\begin{Highlighting}[]
\CommentTok{\# Pivot tract data long on all sex x age columns}
\NormalTok{tract\_long }\OtherTok{\textless{}{-}}\NormalTok{ tract }\SpecialCharTok{|\textgreater{}}
  \FunctionTok{pivot\_longer}\NormalTok{(}
    \AttributeTok{cols =}\NormalTok{ MaleLT5Yrs}\SpecialCharTok{:}\NormalTok{FemaleGE85Yrs,}
    \AttributeTok{names\_to =} \StringTok{"SexAge"}\NormalTok{,}
    \AttributeTok{values\_to =} \StringTok{"n"}
\NormalTok{  ) }\SpecialCharTok{|\textgreater{}}
  \FunctionTok{mutate}\NormalTok{(}
    \AttributeTok{AgeGroup =} \FunctionTok{case\_when}\NormalTok{(}
      \FunctionTok{str\_detect}\NormalTok{(SexAge, }\StringTok{"18to19|20Yrs|21Yrs|22to24|25to29|30to34"}\NormalTok{) }\SpecialCharTok{\textasciitilde{}} \StringTok{"18{-}34"}\NormalTok{,}
      \FunctionTok{str\_detect}\NormalTok{(SexAge, }\StringTok{"35to39|40to44|45to49"}\NormalTok{)                    }\SpecialCharTok{\textasciitilde{}} \StringTok{"35{-}49"}\NormalTok{,}
      \FunctionTok{str\_detect}\NormalTok{(SexAge, }\StringTok{"50to54|55to59|60to61|62to64"}\NormalTok{)             }\SpecialCharTok{\textasciitilde{}} \StringTok{"50{-}64"}\NormalTok{,}
      \FunctionTok{str\_detect}\NormalTok{(SexAge, }\StringTok{"65to66|67to69|70to74|75to79|80to84|GE85"}\NormalTok{) }\SpecialCharTok{\textasciitilde{}} \StringTok{"65+"}\NormalTok{,}
      \AttributeTok{.default =} \ConstantTok{NA\_character\_}
\NormalTok{    )}
\NormalTok{  )}

\CommentTok{\# PSU{-}level (tract) counts by age group}
\NormalTok{psu\_counts }\OtherTok{\textless{}{-}}\NormalTok{ tract\_long }\SpecialCharTok{|\textgreater{}}
  \FunctionTok{filter}\NormalTok{(}\SpecialCharTok{!}\FunctionTok{is.na}\NormalTok{(AgeGroup)) }\SpecialCharTok{|\textgreater{}}
  \FunctionTok{summarise}\NormalTok{(}
    \AttributeTok{n =} \FunctionTok{sum}\NormalTok{(n),}
    \AttributeTok{.by =} \FunctionTok{c}\NormalTok{(Tract, AgeGroup)}
\NormalTok{  ) }\SpecialCharTok{|\textgreater{}}
  \FunctionTok{mutate}\NormalTok{(}\AttributeTok{AgeGroup =} \FunctionTok{factor}\NormalTok{(AgeGroup, }\AttributeTok{levels =} \FunctionTok{c}\NormalTok{(}\StringTok{"18{-}34"}\NormalTok{, }\StringTok{"35{-}49"}\NormalTok{, }\StringTok{"50{-}64"}\NormalTok{, }\StringTok{"65+"}\NormalTok{)))}

\CommentTok{\# SSU{-}level (block group) counts by age group — same pivot applied to bg}
\NormalTok{bg\_long }\OtherTok{\textless{}{-}}\NormalTok{ bg }\SpecialCharTok{|\textgreater{}}
  \FunctionTok{pivot\_longer}\NormalTok{(}
    \AttributeTok{cols =}\NormalTok{ MaleLT5Yrs}\SpecialCharTok{:}\NormalTok{FemaleGE85Yrs,}
    \AttributeTok{names\_to =} \StringTok{"SexAge"}\NormalTok{,}
    \AttributeTok{values\_to =} \StringTok{"n"}
\NormalTok{  ) }\SpecialCharTok{|\textgreater{}}
  \FunctionTok{mutate}\NormalTok{(}
    \AttributeTok{AgeGroup =} \FunctionTok{case\_when}\NormalTok{(}
      \FunctionTok{str\_detect}\NormalTok{(SexAge, }\StringTok{"18to19|20Yrs|21Yrs|22to24|25to29|30to34"}\NormalTok{) }\SpecialCharTok{\textasciitilde{}} \StringTok{"18{-}34"}\NormalTok{,}
      \FunctionTok{str\_detect}\NormalTok{(SexAge, }\StringTok{"35to39|40to44|45to49"}\NormalTok{)                    }\SpecialCharTok{\textasciitilde{}} \StringTok{"35{-}49"}\NormalTok{,}
      \FunctionTok{str\_detect}\NormalTok{(SexAge, }\StringTok{"50to54|55to59|60to61|62to64"}\NormalTok{)             }\SpecialCharTok{\textasciitilde{}} \StringTok{"50{-}64"}\NormalTok{,}
      \FunctionTok{str\_detect}\NormalTok{(SexAge, }\StringTok{"65to66|67to69|70to74|75to79|80to84|GE85"}\NormalTok{) }\SpecialCharTok{\textasciitilde{}} \StringTok{"65+"}\NormalTok{,}
      \AttributeTok{.default =} \ConstantTok{NA\_character\_}
\NormalTok{    )}
\NormalTok{  )}

\NormalTok{ssu\_counts }\OtherTok{\textless{}{-}}\NormalTok{ bg\_long }\SpecialCharTok{|\textgreater{}}
  \FunctionTok{filter}\NormalTok{(}\SpecialCharTok{!}\FunctionTok{is.na}\NormalTok{(AgeGroup)) }\SpecialCharTok{|\textgreater{}}
  \FunctionTok{summarise}\NormalTok{(}
    \AttributeTok{n =} \FunctionTok{sum}\NormalTok{(n),}
    \AttributeTok{.by =} \FunctionTok{c}\NormalTok{(Tract, BlockGroup, AgeGroup)}
\NormalTok{  ) }\SpecialCharTok{|\textgreater{}}
  \FunctionTok{mutate}\NormalTok{(}\AttributeTok{AgeGroup =} \FunctionTok{factor}\NormalTok{(AgeGroup, }\AttributeTok{levels =} \FunctionTok{c}\NormalTok{(}\StringTok{"18{-}34"}\NormalTok{, }\StringTok{"35{-}49"}\NormalTok{, }\StringTok{"50{-}64"}\NormalTok{, }\StringTok{"65+"}\NormalTok{)))}
\end{Highlighting}
\end{Shaded}



  \bibliography{references.bib}



\end{document}
